% cv-skills-section.tex — Skills section component
% Corresponds to SkillsSection.tsx in the web app

% ═════════════════════════════════════════════════════════════════════
% SKILLS SECTION
% Usage: \begin{skillssection}{Title}
%          \skillcategory{\faIcon{code}}{Category}{Skill 1, Skill 2}
%        \end{skillssection}
% ═════════════════════════════════════════════════════════════════════
\newenvironment{skillssection}[1]{%
  \cvsectionheader{#1}%
}{%
  \vspace{0.2cm}%
}

% ─── Single Skill Tag (internal) ─────────────────────────────────────
\newcommand{\skilltag}[1]{%
  \begingroup
  \setlength{\fboxsep}{2.5pt}%
  \colorbox{cv-section-bg}{%
    \kern1pt{\small\color{cv-primary}\rule[-0.5ex]{1.5pt}{2ex}}\kern3pt%
    {\small\color{cv-on-surface}#1}%
    \kern2pt%
  }%
  \endgroup
  \hspace{0.15em}%
}

% ─── Process comma-separated skill list (internal) ───────────────────
\newcommand{\skilltags}[1]{%
  \forcsvlist{\skilltag}{#1}%
}

% ═════════════════════════════════════════════════════════════════════
% SKILL CATEGORY CARD (internal — use \skillrow to render pairs)
% ═════════════════════════════════════════════════════════════════════
\newsavebox{\cvskillcontentbox}
\newcommand{\cvskillcard}[4]{%
  % #1 = width, #2 = min height, #3 = icon+title, #4 = skill tags
  % Measure content
  \savebox{\cvskillcontentbox}{%
    \begin{minipage}{#1-16pt}%
      #3%
      \par\vspace{0.3em}%
      \setlength{\baselineskip}{1.8em}%
      \setlength{\lineskip}{0.3em}%
      \setlength{\lineskiplimit}{0.3em}%
      \skilltags{#4}%
    \end{minipage}%
  }%
  \begin{tikzpicture}
    \node[
      draw=cv-skill-border,
      fill=cv-skill-bg,
      rounded corners=4pt,
      inner sep=8pt,
      text width=#1-16pt,
      minimum height=#2,
      anchor=north west,
    ] (card) {%
      #3%
      \par\vspace{0.3em}%
      \setlength{\baselineskip}{1.8em}%
      \setlength{\lineskip}{0.3em}%
      \setlength{\lineskiplimit}{0.3em}%
      \skilltags{#4}%
    };
  \end{tikzpicture}%
}

% ═════════════════════════════════════════════════════════════════════
% SKILL ROW — renders two cards side by side with equal height
% Usage: \skillrow{icon1}{title1}{skills1}{icon2}{title2}{skills2}
% ═════════════════════════════════════════════════════════════════════
\newlength{\cvskillcolw}
\newlength{\cvskillha}
\newlength{\cvskillhb}
\newlength{\cvskillhmax}
\newsavebox{\cvskillboxa}
\newsavebox{\cvskillboxb}
\newcommand{\skillrow}[6]{%
  \setlength{\cvskillcolw}{0.5\linewidth-0.3em}%
  % Measure card A
  \savebox{\cvskillboxa}{%
    \cvskillcard{\cvskillcolw}{0pt}{{\color{cv-primary}#1}\enspace{\bfseries\color{cv-on-surface}#2}}{#3}%
  }%
  % Measure card B
  \savebox{\cvskillboxb}{%
    \cvskillcard{\cvskillcolw}{0pt}{{\color{cv-primary}#4}\enspace{\bfseries\color{cv-on-surface}#5}}{#6}%
  }%
  % Calculate max height
  \setlength{\cvskillha}{\dimexpr\ht\cvskillboxa+\dp\cvskillboxa\relax}%
  \setlength{\cvskillhb}{\dimexpr\ht\cvskillboxb+\dp\cvskillboxb\relax}%
  \ifdim\cvskillha>\cvskillhb
    \setlength{\cvskillhmax}{\cvskillha}%
  \else
    \setlength{\cvskillhmax}{\cvskillhb}%
  \fi
  % Render both at max height
  \noindent
  \cvskillcard{\cvskillcolw}{\cvskillhmax}{{\color{cv-primary}#1}\enspace{\bfseries\color{cv-on-surface}#2}}{#3}%
  \hfill
  \cvskillcard{\cvskillcolw}{\cvskillhmax}{{\color{cv-primary}#4}\enspace{\bfseries\color{cv-on-surface}#5}}{#6}%
  \par\vspace{0.4em}%
}
